\chapter{图像和多模态处理概述}
\label{chap:vision-overview}

同一段街道录像可以支持许多不同的请求。交通统计关心经过多少辆车，图像检索希望找回拍到某辆车的其他照片，夜景增强希望看清路边招牌，视频编辑则可能要求把车身改成另一种颜色。这些任务都读取视觉信息，但需要保留的证据、允许改变的内容和应当返回的结果并不相同。仅仅选择一个能够接收图像的模型，还没有完整规定要解决的问题。

前面的模型与范式介绍了表示、预测、生成和学习的通用机制。本部分从应用端继续追问：观测怎样进入计算，模型输出怎样还原成对象、区域、时间区间或媒体成果，怎样判断结果满足了最初的请求？图像和视频任务建立空间与时间的基础，多模态模型解释不同信息怎样连接，多模态理解和生成再处理查询、证据及多种输出之间的关系。本章先明确共同对象，随后用任务地图、研究演进与尚未满足的约束说明这些问题如何联系起来。

\section{图像和多模态处理的基本概念}
\label{sec:vision-overview-concepts}

描述一个处理问题，需要同时给出观测、成果和约束。观测决定算法实际上看到了什么，成果决定预测变量是什么，约束则规定哪些误差不能接受。把这三项混在一起，容易把“图像变清晰了”误认为“图中文字恢复正确了”，或把“能输出一个框”误认为“找到了请求中的对象”。

\subsection{处理对象与观测范围}
\label{sec:vision-overview-observations}

一幅数字图像把连续的成像结果记录在离散网格上。\term{像素}（pixel）是网格中的一个位置及其记录值，\term{通道}（channel）则区分该位置记录的不同分量。常见RGB图像有红、绿、蓝三个颜色通道，但通道也可以承载深度、光谱或其他测量。以下把图像写为三阶数组
\begin{equation}
  \bm{X}\in\mathbb{R}^{h\times w\times c},
  \qquad x_{i,j,k}=\text{第$i$行、第$j$列、第$k$通道的值。}
  \label{eq:vision-image-array}
\end{equation}
这里$h,w,c$分别是高度、宽度与通道数。数组字形沿用本书的粗体约定，元素值是标量。输入范围也属于定义的一部分：8位整数的$0$至$255$与归一化后的$0$至$1$不能直接混用。

\term{分辨率}（resolution）在数字图像接口中通常指网格尺寸，例如$1920\times1080$；它决定有多少采样位置，却不能独自说明可辨细节。失焦、高噪声或压缩过重都可能使较大的图像缺少有效信息。镜头、曝光、光照、传感器和后续成像处理影响像素值，同一物体在不同照片中的像素未必相同。把一幅小图插值到更大网格，也没有增加新的实拍观测。

像素坐标只是图上的位置。本部分在没有另行说明时使用横向向右、纵向向下的图像坐标；连续框边界用两个坐标之差计算宽高。若讨论离散像素中心、图像缩放或归一化坐标，会明确给出变换。对象身份用于说明不同位置是否属于同一实体，三维坐标用于说明实体位于哪个空间位置。两辆车可在不同时间经过同一像素，一辆车也可以跨越许多像素，因此这三种描述不能互换。

\term{视频}（video）把画面与时间关联起来，可表示为$\{(\bm{X}_t,\tau_t)\}_{t=1}^{n}$。帧编号$t$表示顺序，时间戳$\tau_t$表示拍摄或播放时刻。固定帧率$f$时，相邻帧间隔为$1/f$；改变抽帧步长会改变模型实际看到的时间间隔。可变帧率或丢帧时，应使用真实时间戳，不能简单把帧编号当秒数。镜头切换会中断局部运动的连续性，相机移动也会让静止物体在图上移动。

文字、图像和声音还可以共同描述一个事件。\term{模态}（modality）指信息的类型或通道；\term{多模态处理}（multimodal processing）利用这些信息之间的联系完成任务。图像可以显示谁在敲门，声音可以提供敲击节奏，文字可以指定应当寻找哪一次敲门。三者的联系需要学习或建立对应，不能由“同时输入”自动得到。同一对象的多张照片增加视角，图像与文字配对增加信息类型；视频在本部分按时序视觉输入展开，页面布局则作为内容的空间组织处理。

\begin{example}[采样之后还能看见什么]
  \label{ex:vision-overview-sampling}
  一段固定帧率为25帧每秒的视频持续12秒，含300帧。若每隔25帧保留一帧，得到的12幅画面间隔1秒。设一辆车只在4.2秒至4.6秒间出现在画面内，而保留帧恰好位于整数秒，则这些输入中完全没有这辆车。后续模型即使正确处理每幅输入，也不能据此证明原视频中没有车辆出现。
\end{example}

这个例子把信息获取和模型计算分开了。错误可能发生在进入模型之前；增加语言推理长度不能补回从未保留的那段观测。长视频、高分辨率页面和细小目标都需要检查同样的问题。

\subsection{处理目标与输出形式}
\label{sec:vision-overview-outputs}

当请求是“图中有没有公交车”时，结果可以是类别判断；请求变为“公交车在哪里”时，还需要空间范围。\term{目标检测}（object detection）通常返回对象的边界框、类别与得分，\term{图像分割}（image segmentation）则为像素指定类别或对象归属。框包含的背景多少、两个同类对象是否分开，都影响结果含义。跟踪进一步把不同时刻的范围关联为同一对象的轨迹。

数量、关节位置和距离属于测量结果。按检测实例计数需要规定怎样处理边界上的对象；从图像估计距离则需要交代相机、尺度和坐标依据。检索返回满足相关性要求的记录及排序。“外观像同类汽车”和“拍到了同一辆汽车”是不同的相关性，不能用同一组正例默认替代。

另一类请求需要返回图像本身。\term{图像复原}（image restoration）依据退化观测估计原内容，例如从带噪照片估计较少噪声的图像；\term{图像增强}（image enhancement）改善指定用途下的呈现，例如调整暗部层次；\term{图像压缩}（image compression）用较少比特表示图像，并规定解码后的允许误差。三者可以组成一条处理链，但分别回答内容估计、呈现和表示成本的问题。

\term{图像生成}（image generation）根据分布或条件创建图像，\term{图像编辑}（image editing）则对已有图像实施指定变化。复原与生成都可能产生原输入中看不到的细节，区别在于任务要求：复原希望估计原场景，创作可以按要求改变场景。一个去模糊模型把招牌上的字改成另一个清晰的字，即使观看效果很好，也没有完成文字保真的复原目标。

图\ref{fig:vision-task-map}把同一街景可能形成的成果并列起来。各分支的主要区别落在输出对象和核验关系上，不意味着必须由完全独立的模型实现。一个共享主干可以连接多个预测头，但仍要分别检查框、掩码、距离和语言描述。

\begin{figure}[htbp]
  \centering
  % 教学任务地图；宽约108mm，采用项目语义配色。
\begin{tikzpicture}[
  font=\figurefont\small,
  task/.style={draw=BookRule,fill=BookPaper,minimum width=30mm,
    minimum height=9mm,align=center,inner sep=2mm},
  link/.style={-{Stealth[length=2mm]},draw=BookMuted,line width=0.8pt},
  note/.style={anchor=west,text=BookMuted,font=\figurefont\footnotesize}
]
  \node[task,fill=FigureInput!10,draw=FigureInput,minimum width=20mm] (input) at (0,0) {街景观测};
  \node[task] (recognition) at (4,2.6) {类别、属性};
  \node[task] (region) at (4,1.3) {框、掩码、轨迹};
  \node[task] (measure) at (4,0) {数量、距离、排序};
  \node[task] (restore) at (4,-1.3) {复原、压缩结果};
  \node[task] (create) at (4,-2.6) {编辑、描述、配声};
  \draw[draw=BookMuted,line width=0.8pt] (input.east) -- (1.8,0);
  \draw[draw=BookMuted,line width=0.8pt] (1.8,-2.6) -- (1.8,2.6);
  \foreach \name/\y in {recognition/2.6,region/1.3,measure/0,restore/-1.3,create/-2.6}{
    \draw[link] (1.8,\y) -- (\name.west);
  }
  \node[note] at (5.8,2.6) {对象与语义};
  \node[note] at (5.8,1.3) {范围与身份};
  \node[note] at (5.8,0) {尺度与相关性};
  \node[note] at (5.8,-1.3) {保真与比特};
  \node[note] at (5.8,-2.6) {条件与对应};
\end{tikzpicture}

  \caption{同一街景观测可以服务不同任务。各行表示一种成果及其主要核验关系，均为教学构造，不代表某个模型的实测输出。图像、视频、文字和声音可以分别进入相应分支；输入相同不意味着监督或评价相同。}
  \label{fig:vision-task-map}
\end{figure}

\subsection{输入依据与结果约束}
\label{sec:vision-overview-evidence}

处理问题中的信息来源有不同职责。照片和录音是观测，人工类别、框和掩码是训练或评价标注，查询说明希望找什么，创作条件说明希望生成什么，相机参数与曝光记录则提供辅助测量。一句“红色汽车正在转弯”若是检索查询，不能同时被当作视频中确有这一事件的证据。

配对信息也可能不完整。网页文字常只描述图片的一部分，旁白可能先讲下一步操作，字幕可能来自画外人物。需要把图像编号、区域、时间与来源保留下来，再检查是否对应同一实体或事件。模型得分度量的是给定训练与接口下的匹配程度，不是对配对事实的无条件证明。

约束必须落实到可以检查的量。车辆计数可以规定在某时间段越过指定线的个体数；文字读取可以要求完整字符串一致；压缩可以给定每像素比特数；交互分割可以限定用户点击次数。分辨率、允许等待的时间和候选选择预算会改变可用信息，也必须进入比较协议。

数据划分应对应实际使用中的独立单位。如果相邻视频帧随机分入训练与测试，两边可能含有几乎相同的背景和对象，测得的结果就不能充分代表面对新视频的表现。医学图像常需按患者或采集中心划分，工业图像可能需要按批次与设备划分，文档可能需要按来源与模板划分。选用哪一种划分取决于希望估计的未来变化。

ImageNet-C用规定的合成退化检查分类器面对常见扰动时的变化，WILDS则把真实使用中的地点、设备、时间等分布差异纳入数据与协议。二者回答的证据问题不同：在合成噪声上稳定，并不能直接推出跨医院或跨地域也稳定\scite{hendrycks2019corruptions,koh2021wilds}

\section{图像和多模态处理的研究内容}
\label{sec:vision-overview-tasks}

本部分按主要成果及新增约束安排应用。图像任务建立单幅观测的输出要求，视频任务增加跨帧对应和信息到达条件，多模态任务增加查询与证据、不同输入及不同输出之间的联系。多模态模型的共用计算集中在\ref{chap:vision-multimodal-models}章，各应用章在需要时回指，避免每次使用文字条件都重讲同一套模型。

\subsection{图像应用的研究内容}
\label{sec:vision-overview-image-tasks}

在工业质检中，“这个零件是否异常”和“异常位于哪里”是两个相互联系的任务。正常样本的局部特征可以给出偏离程度，区域结果还需与具体缺陷对应。若照片本身严重模糊，复原可以改善可见性，但生成出的划痕不能作为真实缺陷的新证据。图像识别、区域判断与复原应在相同对象上保留来源关系。

照片整理提供另一组联系。场景类别可以先缩小候选，再由全局图像表示找相似照片，局部对应与几何验证进一步检查是否拍到了同一建筑。\ref{chap:vision-image-recognition}章据此讨论语义、范围、测量与匹配。\ref{chap:vision-image-restoration}章讨论复原、增强和压缩；\ref{chap:vision-image-generation}章解释创建、转换、编辑及观测场景的新视角呈现。它们共享图像表示，但对细节变化的容许程度不同。

\subsection{视频应用的研究内容}
\label{sec:vision-overview-video-tasks}

看到每一帧中的车辆还不足以统计车辆经过路口的数量。统计需要区分同一辆车的重复出现，跟踪需要在遮挡前后维持身份，转弯判断又要把轨迹与道路及时间范围联系起来。\ref{chap:vision-video-recognition}章分别讨论这些输出。逐帧检测正确但身份交换，是跟踪错误；身份正确但事件边界偏移，则是活动定位错误。

邻帧也可以提供单幅图像缺少的细节。某帧的路牌模糊，另一帧可能清晰；只有把两帧中同一位置的内容对应起来，额外观测才可能帮助恢复。\ref{chap:vision-video-restoration}章讨论跨帧选择、对齐、传播、超分辨率、插帧和编码。更多帧同时带来等待、存储与错误传播，双向读取前后帧的算法还不能直接用于不允许等待的在线处理。

\subsection{多模态应用的研究内容}
\label{sec:vision-overview-multimodal-tasks}

“找到车辆转弯的片段”把文字查询与时间证据联系起来；“为什么画面中有人回头”可能还需要听见声音，但可观察到先后关系并不等于识别了因果机制。文档和图表增加了字符串、阅读顺序、行列与单位的对应。\ref{chap:vision-multimodal-understanding}章从检索、定位、读取和结构恢复推进到描述与问答，始终保留答案所依赖的来源。

“生成同一辆车转弯并带发动机声音的片段”则改变了问题性质。这里的车、运动和声音是创作条件，不是待解释的实拍证据。\ref{chap:vision-multimodal-generation}章讨论视频、配声、联合音视频及空间成果的制作。固定已有视频再生成声音，与画面和声音共同生成，具有不同的可修改范围；单幅文生图虽然也是跨模态任务，其完整应用流程集中在图像生成章。

\section{图像和多模态处理的研究历史}
\label{sec:vision-overview-history}

视觉研究没有一条覆盖所有任务的单线发展史。分类、几何、复原、编码与生成长期并行，也不断借用彼此的表示和计算。以下按应用要求的扩展组织几个转折，阶段相互重叠，代表工作只用于解释新增问题。

\subsection{人工设计方法与观测约束}
\label{sec:vision-overview-classical-history}

在直接从大量样本学习整套系统之前，滤波、退化建模和几何约束已为视觉处理提供了可计算依据。已知模糊过程时，可以研究哪些频率受到抑制以及怎样限制反演时的噪声放大；多幅照片有可靠对应时，可以利用相机投影恢复空间结构。观测模型说明信息怎样进入图像，也指出无法从输入唯一确定的部分。

SIFT通过尺度空间中的局部极值、方向分配和局部描述，使图像中部分可重复辨认的位置能够被匹配。它为跨尺度与旋转变化下的局部对应提供了代表路线；这里的适应范围不等于任意视角、遮挡或材质变化下都能可靠匹配\scite{vision-lowe2004}

这类方法没有因深度网络出现而失去全部作用。几何一致性仍可检查学习式匹配结果，退化与采样关系仍可约束复原，码流解码仍须遵守编码端与解码端的信息一致性。后来变化较大的是哪些表示与计算由数据学习，物理观测的限制并未随之消失。

\subsection{学习式视觉应用的扩展}
\label{sec:vision-overview-learned-history}

AlexNet在大规模图像分类中的应用展示了联合学习视觉表示和类别预测的路线。它说明整图类别任务如何利用多层特征，但分类输出本身不包含对象范围\scite{krizhevsky2012}

全卷积网络把预测组织为与空间位置对应的类别图，使分类网络中的表示进入密集输出。检测需要进一步返回多个对象，实例分割需要区分同类个体，全景分割则共同组织可数对象与背景区域。这些变化增加的是输出结构与监督职责，不能只用“网络变深了”解释\scite{vision-1411-4038}

与此同时，视频中的时空表示和对象关联、低照度成像、图像转换及学习式压缩也在发展。它们分别处理时间证据、成像条件、输入输出映射和有限比特表示。一个网络能够在多种任务中使用，不代表同一检查点在没有对应监督或条件接口时就具备所有能力。

\subsection{可复用能力与开放条件}
\label{sec:vision-overview-open-history}

约2020至2023年的若干代表路线进一步改变了应用接口。DETR把对象检测组织为无序集合与匹配问题；CLIP把文字描述接到可比较的图像表示，使用户可以在使用时指定候选类别；生成式视觉语言模型则把视觉信息接入语言预测，支持描述和问答。CLIP的表示匹配与语言模型的句子生成承担不同计算职责\scite{radford2021}

用户还可以用点或框选择待分割对象，用空间条件约束生成布局，用少量参考图像指定创作主体。提示扩展了交互方式，也引入了目标歧义：落在车门上的一点可以指车门，也可以指整辆车。开放词汇和通用接口扩大了候选范围，结果仍须依照具体目标核验。

新视角合成形成了另一条支线。NeRF以及后来的三维高斯表示，依据已拍视图拟合可渲染场景，再计算新相机下的图像。这个任务同时利用观测约束与学习或优化，但照片渲染准确并不自动证明表面几何适合测量。相关机制在\ref{sec:vision-generation-novel-view}节集中展开。

\subsection{连续、多源与交互应用}
\label{sec:vision-overview-continuous-history}

较新的应用进一步要求系统处理连续观察和多种成果。SAM 2把提示分割延伸到有历史记忆的视频，MovieChat研究长视频信息的压缩与取回，VGGT联合预测多视图几何，Qwen2.5-Omni把音画输入与文字、语音响应联系起来。这些工作分别扩展了对象延续、历史证据、空间输出与响应接口，并不具有同一项成功标准；原论文与固定版本在后续对应章节给出。

生成系统也从单幅图片扩展到视频、空间内容与音视频。共享模型和多阶段制作同时存在：一种路线共享部分表示与主干，另一种路线让专用模型通过明确接口传递结果。仅凭一个产品入口可以接收多种输入，无法判断内部是否共同训练、是否共享参数，更无法推出各项输出之间已保持一致。

\section{当前前沿与研究方向}
\label{sec:vision-overview-frontiers}

研究方向可以从尚未满足的应用约束中找到。以下讨论截至2026年10月5日文献地图收录的代表问题，近期预印本用于说明可检验的研究路线。具体模型的局部结果只支持其数据、版本与观察预算下的判断。

\subsection{高分辨率与长时间中的证据保存}
\label{sec:vision-overview-evidence-frontier}

整幅页面压缩到有限网格后，小字可能消失；长视频均匀采样后，短暂动作可能没有进入输入。问题不只是容纳更多词元，还包括怎样选择局部区域、何时再次读取原始内容，以及怎样保存以后可能需要的细节。压缩历史与完整历史在信息上不同，语言摘要不能无条件替代原观测。

VideoZeroBench把答案、时间区间和空间框分层核验，用于检查“答对了”与“找对证据”是否一致。这样的任务提示我们：应同时评价证据获取与结果形成，不能仅凭最终答案反推中间观察一定正确。\ref{sec:vision-understanding-evaluation}节将展开相应诊断。

\subsection{空间、时间与主体状态的一致性}
\label{sec:vision-overview-consistency-frontier}

相机、深度、点图与轨迹若描述同一场景，应在共同坐标及尺度约定下相容。一个输出把车辆放在路灯前方，另一个输出却给出相反的深度关系，即使各自看起来合理，组合结果也不能直接使用。VGGT提供联合几何预测的路线，空间问答基准则检查模型怎样使用这些关系；几何预测和语言答案需要分别评价。

时间上的一致性也有不同对象。车辆遮挡后再次出现，需要核验身份；动作反向播放，需要核验事件顺序；生成视频中的物体离开画面后重现，需要核验主体与状态保持。这些问题都涉及连续性，却不能由一个通用的“稳定分数”替代。

\subsection{不同模态对同一事件的对应}
\label{sec:vision-overview-alignment-frontier}

声音和画面可以互补，也可能冲突。画面中的人嘴唇在动，不表示录音一定来自此人；字幕出现一个名称，也不证明该对象可见。模型需要保留谁提供了什么信息、发生在何时，再决定如何融合。

多源模型提供了联合计算的接口，但评价还要检查删除、替换或错配某种模态后的变化。若问题单凭字幕就能回答，原本的准确率不能独自证明视觉模块有效。模态消融用于定位贡献，其结论仍受替换方式和分布变化影响，不能把一次扰动结果当作完整因果证明。

\subsection{生成内容的控制与跨成果协调}
\label{sec:vision-overview-generation-frontier}

“一只红色玩具狗在蓝杯左侧”同时要求对象存在、数量正确、属性归属和空间关系成立。局部修改还要求未指定对象保持，视频增加持续身份和动作，声音增加声源与时间对应。全局图文相似度只能覆盖其中一部分。

Show-o2研究理解与生成能力的共享，LTX-2研究联合音视频生成。它们为不同输出之间的信息交换提供了机制，却不能由自然的画面或同步的短片段直接推出长叙事、物理接触及全局状态都正确。后续生成章将把条件遵循、内容保持和制作预算分别记录。

\subsection{持续观察与有限预算下的应用}
\label{sec:vision-overview-streaming-frontier}

在线系统只能使用已经到达的信息。设时刻$\tau$能够使用的观测集合为
\begin{equation}
  \mathcal{O}(\tau)=\{(\bm{X}_t,\tau_t):\tau_t\leq\tau\}\text{。}
  \label{eq:vision-causal-observations}
\end{equation}
若算法必须等待$\delta$秒后才能形成输出，它在时刻$\tau+\delta$可以使用更晚的观测，但这段等待应计入响应延迟。图\ref{fig:vision-observation-window}区分了这一信息条件。循环网络、缓存或长记忆描述的是计算组织，是否因果则取决于真正读取了哪些时刻。

\begin{figure}[htbp]
  \centering
  % 时间坐标只作范围示意，不表示某个实验的采样频率。
\begin{tikzpicture}[x=1cm,y=1cm,font=\figurefont\small]
  \fill[FigureInput!10] (0,0.15) rectangle (4.6,1.3);
  \fill[FigureModel!10] (4.6,0.15) rectangle (7,1.3);
  \fill[BookPaper] (7,0.15) rectangle (10.2,1.3);
  \draw[-{Stealth[length=2mm]},BookInk,line width=0.8pt] (0,0) -- (10.5,0)
    node[below left] {观测时间};
  \foreach \x in {0.4,1.4,2.4,3.4,4.4,5.4,6.4,7.4,8.4,9.4}{
    \fill[BookInk] (\x,0) circle (1.5pt);
  }
  \draw[BookInk,line width=0.9pt] (4.6,-0.25) -- (4.6,1.5);
  \draw[BookMuted,dashed,line width=0.8pt] (7,-0.25) -- (7,1.5);
  \node[below] at (4.6,-0.25) {$\tau$};
  \node[below] at (7,-0.25) {$\tau+\delta$};
  \node[align=center] at (2.3,0.75) {已到达历史\\即时可用};
  \node[align=center] at (5.8,0.75) {等待期间\\新增观测};
  \node[align=center] at (8.6,0.75) {其余未来\\离线才可用};
\end{tikzpicture}

  \caption{同一查询在不同信息条件下可读取的帧。圆点表示实际观测，竖线是查询发起时刻$\tau$；在线即时处理只读取左侧，允许等待$\delta$时才可加入中间区间，离线处理还可使用右侧的给定未来。颜色之外用文字和线型区分范围。}
  \label{fig:vision-observation-window}
\end{figure}

Flash-VStream与TempoGround分别涉及历史信息取回、流式目标对应和状态维护。相关研究需要同时记录帧率、历史容量、观察范围与响应时间。存储更多摘要可能保存较长经历，却仍遗漏小目标；处理更多区域可能取得细节，却增加等待。应用选择应从请求所需证据出发，而不是把输入越多、模型越复杂直接当作任务完成。

\section{本章小结}
\label{sec:vision-overview-summary}

街景上的分类、定位、增强与编辑之所以需要不同解法，是因为它们要求不同的成果，也容许不同的内容变化。明确观测、输出和约束之后，才能判断模型需要保留什么表示、取得什么监督，以及怎样恢复和核验结果。图像建立空间关系，视频增加身份、时间与信息到达条件，多模态处理再连接查询、不同来源及不同输出。

研究范围从专门任务扩展到开放、连续和多源应用，并没有取消证据边界。采样可能遗漏内容，预测可能依赖先验，生成可能补出没有观测支持的细节。后续章节据此逐项展开方法：既追踪输入怎样进入计算，也检查结果是否仍指向正确的对象、区域、时刻和用途。
